\section{From the pinned comparison to magnetization}\label{sec:pressure}

We first combine the two residual estimates to prove the uniform
pinned comparison.  We then express a positive field as a mixture
of deterministic pin sets.  Integrating the finite-volume
magnetization bounds in the field will preserve the required order:
volume tends to infinity before the field tends to zero.

\subsection{Completion of the capacity comparison}

\begin{proof}[Proof of Proposition~\ref{prop:comparison}]
Fix the cube, pins, and starting slot in the proposition, and expose
the cycles under $\PP_P$, $P=E\cup\{i\}$.
Use the trial vector $h$ and constant $c$ from
\eqref{eq:trialvector} and \eqref{eq:freepotential}.
On the good event $\mathcal G$ of Proposition~\ref{prop:good},
Proposition~\ref{prop:dualenergy}, applied with a fixed supremum
bound $2$, gives a measurable $X$ such that
\[
 \abs{\ip{\mu_d}{z}}\le X\sqrt{e_z}+\varepsilon_\beta,
 \qquad d=+,-,
\]
for vectors vanishing on killed cut slots with
$\norm{z}_\infty\le2$, where $\varepsilon_\beta$ is rapidly small
and
\[
 \EE_P[\ind_{\mathcal G}X^2]=o(\beta^{-5/2}).
\]
If the two directions initially give different coefficients, take
their maximum; this preserves the displayed expectation bound.
The source identities \eqref{eq:residualdefinition} identify these
residuals with those in \eqref{eq:capacityresiduals}.
Since $0\le h\le1$, the averaged capacity criterion
\eqref{eq:averagedcapacity} applies.  Proposition~\ref{prop:scalar}
and the probability estimate in Proposition~\ref{prop:good} give
\begin{equation}\label{eq:returncomparison}
 r:=\EE_Pq_i=1-c+o(\beta^{-5/2}).
\end{equation}
Every error here is uniform over the cube and the allowed pin sets.
In particular, the good-event estimate is used only through its
probability on the complement; there $q_i\in[0,1]$ suffices.

By Lemma~\ref{lem:freebounds}, $v(x)=O(\beta^{-3/2})$, so
\[
 1-c=\frac{v(x)}{1+v(x)}=v(x)+O(\beta^{-3}).
\]
The exact pin identity \eqref{eq:pinreturn} now yields
\begin{equation}\label{eq:slotcomparison}
 d_E(i)=\frac{r}{1+r}
       =v(x)+o(\beta^{-5/2})
       =\frac1\ell\sum_{n\ge1}p_{n\beta}(0)
          +o(\beta^{-5/2}).
\end{equation}
Indeed $r=O(\beta^{-3/2})$, so the second rational transformation
also costs only $O(\beta^{-3})$.
By \eqref{eq:freeidealdensity}, the last sum is
$n_S^{\rm sw}(\beta)$.  The uniform little-oh
in \eqref{eq:slotcomparison} is equivalent to the asserted function
$\varepsilon_S(\beta)\to0$ after choosing a sufficiently large
$\beta_0(S)$.  Equations \eqref{eq:returncomparison} and
\eqref{eq:freeidealdensity}, together with
$1-c=v(x)+O(\beta^{-3})$, also give
$r=\ell^{-1}n_S^{\rm sw}(\beta)+o(\beta^{-5/2})$, as asserted.
\end{proof}

\subsection{Boundary pins and field pins}

Write $V_N=|\Lambda_N|$, use
$\partial\Lambda_N=\partial_{\mathrm{in}}\Lambda_N$ for the inner
boundary, and let $q_{\Lambda_N}$ be the number of slot edges.
For a picture $\omega$, let $\mathcal C(\omega)$ be
its cycles and $a_C$ the number of time-zero slots on $C$.
For colors $\sigma_C\in\{-1,1\}$, define
\[
 M(\omega,\sigma)=\frac12\sum_Ca_C\sigma_C,
 \qquad
 N_\downarrow(\omega,\sigma)=\sum_{C:\,\sigma_C=-1}a_C
                              =SV_N-M(\omega,\sigma).
\]
Let $\chi_D(\omega,\sigma)$ indicate that every cycle meeting $D$
is up.  The spin-loop representation in
Proposition~\ref{prop:loopinputs} identifies the ordinary partition
function with
\[
 Z_N(t)=e^{\beta q_{\Lambda_N}/4}
       \EE_{\rm b}\sum_\sigma e^{\beta tM(\omega,\sigma)},
 \qquad p_N(t)=\frac{\log Z_N(t)}{\beta V_N}.
\]
Its boundary-pinned counterpart and pressure are
\begin{equation}\label{eq:pinnedpressure}
 Z_N^+(t)=e^{\beta q_{\Lambda_N}/4}
       \EE_{\rm b}\sum_\sigma
          \chi_{D_0}(\omega,\sigma)e^{\beta tM(\omega,\sigma)},
 \qquad
 p_N^+(t)=\frac{\log Z_N^+(t)}{\beta V_N}.
\end{equation}

We use the surface estimate of
\cite[Proposition~7.1]{BlochLaw}.  Its argument is short and records why
all-times boundary pins do not change the pressure.

\begin{lemma}[Surface cost of boundary pins]\label{lem:surfacecost}
For every finite $\beta>0$, every $t\ge0$, and every finite cube,
\begin{equation}\label{eq:surfacecost}
 0\le p_N(t)-p_N^+(t)
 \le (\log2)\left(\frac\ell\beta+24\ell^2\right)
        \frac{|\partial\Lambda_N|}{V_N}.
\end{equation}
In particular $p_N^+(t)\to p_{S,\beta}(t)$ for each fixed
$\beta>0$ and $t\ge0$.
\end{lemma}

\begin{proof}
After extracting the all-up factor, the field-tilted picture weight is
\[
 W_t(\omega)=\prod_C(1+e^{-\beta t a_C}).
\]
Let $\EE_t$ denote expectation under this tilted picture law.
Given the picture, each cycle is up independently with probability
$(1+e^{-\beta t a_C})^{-1}\ge1/2$.  If $b_0$ counts cycles meeting
$D_0$, Jensen's inequality therefore gives
\[
 1\ge\frac{Z_N^+(t)}{Z_N(t)}
   \ge\EE_t2^{-b_0}
   \ge\exp[-(\log2)\EE_t b_0].
\]
Let $J_\partial$ count marks on slot edges incident to a boundary
site, and let $q_\partial$ be the number of those edges.
Every boundary-meeting cycle either has a boundary slot at time zero
or passes through one of these marks.  A mark belongs to at most two
cycles, so
\[
 b_0\le\ell|\partial\Lambda_N|+2J_\partial,
 \qquad q_\partial\le6\ell^2|\partial\Lambda_N|.
\]
Insertion of one mark joins two cycles or splits one.  Each factor
of $W_t$ lies in $[1,2]$, and at most two new factors appear, whence
$W_t(\omega\cup\{u\})\le4W_t(\omega)$ for almost every inserted
mark $u$.  The Poisson insertion formula at rate $1/2$ implies
\[
 \EE_t J_\partial\le2\beta q_\partial
               \le12\beta\ell^2|\partial\Lambda_N|.
\]
The preceding inequalities, divided by $\beta V_N$, prove
\eqref{eq:surfacecost}.  The pressure limit follows because the
boundary-to-volume ratio tends to zero.
\end{proof}

For $t\ge0$ put $b_\beta(t)=1-e^{-\beta t}$, and let $\nu_b$
be the law that selects each time-zero slot independently with
probability $b=b_\beta(t)$.  If $A$ is the selected set and
$\operatorname{Up}(\omega,\sigma)$ is the set of up slots at that
cut, then
\begin{equation}\label{eq:fieldpinidentity}
 \sum_A\nu_b(A)\ind_{\{A\subset\operatorname{Up}\}}
   =(1-b)^{N_\downarrow}=e^{-\beta tN_\downarrow}.
\end{equation}
Consequently the field law in \eqref{eq:pinnedpressure} is the
picture/color marginal of the joint law
\begin{equation}\label{eq:fieldjointlaw}
 \widehat\PP_{N,t}(d\omega,d\sigma,A)
  \ \propto\ \PP_{\rm b}(d\omega)\,\nu_b(A)
                         \chi_{D_0\cup A}(\omega,\sigma),
\end{equation}
where the color variable uses counting measure.
Indeed the common factor extracted from its normalizer is
$e^{\beta q_{\Lambda_N}/4+\beta tSV_N}$.
Conditional on $A$, this is exactly the zero-field color law with
the deterministic pins $E=D_0\cup A$.

The marginal law of $A$ is generally not Bernoulli, because its
normalization is weighted by the pinned partition function.
Conditioning instead on $(\omega,\sigma)$ makes selections on up
slots independent Bernoulli choices of parameter $b$, and forbids
selections on down slots.  Thus every specified slot $j$ satisfies
\begin{equation}\label{eq:fieldpinbound}
 \widehat\PP_{N,t}(j\in A)
   =b_\beta(t)\widehat\PP_{N,t}(j\text{ up})\le b_\beta(t).
\end{equation}
This is the only property of the posterior pin law needed below.

\subsection{Thermodynamic and zero-field limits}

\begin{proof}[Proof of Theorem~\ref{thm:main}]
Fix $S$ and a sufficiently large $\beta$.  The finite-volume
magnetization is bounded, so differentiation under the sum in
\eqref{eq:pinnedpressure} is justified.  Equations
\eqref{eq:fieldjointlaw} and \eqref{eq:fieldpinidentity} give
\begin{equation}\label{eq:pressuredeficit}
 S-(p_N^+)'(t)
  =\frac1{V_N}\widehat\EE_{N,t}N_\downarrow
 =\frac1{V_N}\sum_{j\text{ at time }0}
                         \widehat\EE_{N,t}d_{D_0\cup A}(j).
\end{equation}
Here $d_E(j)$ means the down probability also for pinned slots, where
it is zero.
Let $\rho_N$ be the fraction of sites $x$ for which
$B_{50R_*}(x)$ is interior to $\Lambda_N$.
For such a site let $\mathcal A_x$ be the event that no selected
pin lies above that ball.  A union bound using
\eqref{eq:fieldpinbound} gives
\[
 \widehat\PP_{N,t}(\mathcal A_x^c)
 \le C_{S,\beta}\,b_\beta(t),
\]
where $C_{S,\beta}$ is finite and independent of $N$.
On $\mathcal A_x$, Proposition~\ref{prop:comparison} applies to
each of the $\ell$ slots at $x$, for every conditional pin set $A$.
On its complement, and at sites outside the interior region, use
the boundedness of down probabilities.  Summing the slot estimates
in \eqref{eq:pressuredeficit} gives
\begin{equation}\label{eq:derivativecomparison}
 \abs{S-(p_N^+)'(t)-n_S^{\rm sw}(\beta)}
 \le \ell\varepsilon_S(\beta)\beta^{-5/2}
       +C_S(1-\rho_N)+C_{S,\beta}b_\beta(t).
\end{equation}
Here $n_S^{\rm sw}(\beta)$ is bounded for the large values of $\beta$
under consideration, so its contribution on exceptional
sites is included in $C_S$.

Integrate \eqref{eq:derivativecomparison} over $t\in[0,h]$, for
$h>0$, and divide by $h$.  Now take $N\to\infty$, keeping
$\beta$ and $h$ fixed.  The radius is fixed in this limit, so
$\rho_N\to1$, and Lemma~\ref{lem:surfacecost} gives convergence
of $p_N^+$ at both field endpoints.  We obtain
\[
 \abs{S-\frac{p_{S,\beta}(h)-p_{S,\beta}(0)}h
                    -n_S^{\rm sw}(\beta)}
 \le \ell\varepsilon_S(\beta)\beta^{-5/2}
       +\frac{C_{S,\beta}}h\int_0^h b_\beta(t)\,dt.
\]
For this fixed $\beta$, the last term tends to zero as
$h\downarrow0$.  Convexity and the uniform magnetization bound
give the right derivative of the limiting pressure.  Therefore
\[
 \abs{S-m_S(\beta)-n_S^{\rm sw}(\beta)}
 \le\ell\varepsilon_S(\beta)\beta^{-5/2}.
\]
Finally let $\beta\to\infty$.  Since $S$ is fixed, multiplication
by $(\beta S)^{5/2}$ proves the theorem.
\end{proof}

\subsection{The explicit lattice correction}

The full ideal density has a two-term expansion with a remainder
small enough to identify the coefficient in
Corollary~\ref{cor:expansion}.

\begin{proof}[Proof of Corollary~\ref{cor:expansion}]
Put $t=\beta S$.  Independence of the three coordinates gives
\[
 p_\beta(0)=I(t)^3,\qquad
 I(t)=\frac1{2\pi}\int_{-\pi}^{\pi}
                         e^{-2t(1-\cos k)}\,dk.
\]
We claim, uniformly for $t\ge1$, that
\begin{equation}\label{eq:laplaceexpansion}
 I(t)=\frac1{\sqrt{4\pi t}}
          \left(1+\frac1{16t}+O(t^{-2})\right).
\end{equation}
To see this, set $u=\sqrt t\,k$ and first restrict to
$|u|\le t^{1/10}$.  Taylor's formula gives
\[
 2t(1-\cos(u/\sqrt t))
       =u^2-\frac{u^4}{12t}+O(u^6/t^2),
\]
and hence, on this interval,
\[
 e^{-2t(1-\cos(u/\sqrt t))}
 =e^{-u^2}\left(1+\frac{u^4}{12t}
                   +O\!\left(\frac{u^6+u^8}{t^2}\right)\right).
\]
The bound $2(1-\cos k)\ge c k^2$ controls the omitted interval by
$O(e^{-c t^{1/5}})$ after the same rescaling.  Extending the
Gaussian integrals to $\R$ has an error of this size as well.
Since
\[
 \int_\R e^{-u^2}\,du=\sqrt\pi,
 \qquad
 \int_\R u^4e^{-u^2}\,du=\frac34\sqrt\pi,
\]
the formula \eqref{eq:laplaceexpansion} follows, with a constant
that can be enlarged to cover all $t\ge1$.
Cubing gives
\[
 p_\beta(0)
   =(4\pi t)^{-3/2}\left(1+\frac3{16t}+O(t^{-2})\right).
\]
Apply this uniform estimate at $nt$, $n\ge1$, and sum.
The remainder is bounded by
$Ct^{-7/2}\sum_{n\ge1}n^{-7/2}$, so
\begin{equation}\label{eq:idealdensityexpansion}
 n_S^{\rm sw}(\beta)
  =\frac{\zeta(3/2)}{8\pi^{3/2}}(\beta S)^{-3/2}
   +\frac{3\zeta(5/2)}{128\pi^{3/2}}(\beta S)^{-5/2}
   +O((\beta S)^{-7/2}).
\end{equation}
Theorem~\ref{thm:main} now gives the asserted expansion of
$S-m_S(\beta)$.
\end{proof}
